Das Zahlenverhältnis zwischen bestehenden Ehen und Witwen, das in \textsc{Kritter}s Gutachten für die Professorenwitwenkasse mit \(2:1\) vorausgesetzt wird, kann bei Gesellschaften aus verschiedenen Lebenskreisen sehr verschieden sein. Vergleichen wir zum Beispiel eine Witwenkassengesellschaft wie die unsrige mit der Gesamtheit aller bestehenden Ehen und Witwen in einem ganzen Land zunächst nur hinsichtlich des allgemeinen, [S. 140 oben] angegebenen Bestimmungsgrunds, so sieht man leicht, dass in jener verhältnismäßig mehr Witwen auf eine bestimmte Zahl von Ehen gerechnet werden müssen als in dieser. Bei der großen Masse der Einwohner eines Landes finden die Eheschließungen durchschnittlich in einem früheren Alter statt, und der durchschnittliche Altersunterschied zwischen Mann und Frau ist geringer als bei Universitätsprofessoren\footnote{Aus unserer Witwenkasse liegen mir Altersdaten nur für 62 Ehepaare vor; danach beträgt der durchschnittliche Unterschied 9 Jahre. Für die Gesamtheit der Einwohner eines ganzen Landes wird man schwerlich auch nur einen halb so großen durchschnittlichen Unterschied annehmen können.}. Dazu kommt, dass manche Mitglieder schon verheiratet in unsere Witwenkasse eintreten, während in den Listen eines ganzen Staates sämtliche Ehepaare gleich von ihrer Heirat an mitgerechnet werden. Die Wirkung dieser beiden Ursachen ist jedoch nur gering im Vergleich zum Einfluss eines anderen Umstands, der oben S. [140 ff.] bei der abstrakten Behandlung noch beiseitegelassen wurde: Witwen können nicht allein durch Tod, sondern auch durch Wiederverheiratung aus dem Bestand ausscheiden. In einer Witwenkasse, in der eine Witwe bei Wiederverheiratung jeden weiteren Anspruch auf die Pension verliert, kommen Wiederverheiratungen gewöhnlich selten vor. Insbesondere zählt unsere Witwenkasse in der gesamten Zeit ihres Bestehens nur einen einzigen solchen Fall. Für ein ganzes Land hingegen nimmt man aufgrund der Erfahrung an, dass von der Gesamtzahl der Witwen jährlich etwa der dreißigste Teil durch Wiederverheiratung ausscheidet. Dadurch wird das Verhältnis der bestehenden Ehen zu den Witwen wesentlich verändert. Schließlich hat auch in einem Land, dessen Bevölkerung schon seit längerer Zeit zunimmt, diese Zunahme einen wesentlichen Einfluss auf das Verhältnis der \textit{gleichzeitig bestehenden} Ehen und Witwen: Auf eine Witwe kommen dann mehr bestehende Ehen, als ohne diesen Umstand vorhanden wären.

Als Erfahrungssatz wird gewöhnlich aufgestellt, dass für ein ganzes Land vier bestehende Ehen auf eine Witwe zu rechnen seien. Die von \textsc{Quetelet} für Belgien angegebenen Zahlen stimmen mit diesem Verhältnis fast genau überein. Für das ganze Königreich Hannover ergeben die Zählungen von 1833—1842 eine etwas kleinere Zahl, nämlich 3,74 bestehende Ehen auf eine Witwe. Unterscheidet man aber die einzelnen Landesteile, so zeigen sich sehr große Unterschiede; zum Beispiel kommen
\SourceLeaf{150}
\begin{center}
\begin{tabular}{l@{\;}r}
in der Landdrostei Stade .\,.\,.\,.\,.\,.\,.\,.\,.\,.\,.\,. & 4,21\\
in der Landdrostei Aurich .\,.\,.\,.\,.\,.\,.\,.\,.\,.\,. & 3,28\\
in der Berghauptmannschaft Clausthal .\,.\,.\,. & 2,56
\end{tabular}
\end{center}
\noindent bestehende Ehen auf eine Witwe.

{\rightskip=0pt plus 1em
In der zweiten Hälfte des vorigen Jahrhunderts hat man manche Witwenkassen zugrunde gehen sehen, weil sie auf der nach den obigen Bemerkungen ganz falschen Voraussetzung beruhten\footnote{Manche allerdings auch infolge \textit{noch} gröberer Fehler.}, dass das Verhältnis der bestehenden Ehen zu den Witwen, das sich aus Zählungen für ein ganzes Land ergibt, auch für Witwenkassen als maßgebend betrachtet werden dürfe. \textsc{Kritter} war einer von denen, die nicht müde wurden, solche falschen Grundsätze zu bekämpfen. Wo er aber in seinen Schriften bestimmte eigene Angaben macht, heißt es \textit{entweder} nur, dass man \textit{zum allerwenigsten} eine Witwe auf zwei Ehen rechnen müsse (zum Beispiel \textsc{Kritter}s Vorstellung des bisherigen Erfolgs u. s. w. S. 22); \textit{oder} er verwahrt sich ausdrücklich (Prufung eines Aufsatzes in der Berliner Monatsschrift S. 28) dagegen, das Verhältnis \(2:1\) anders als unter der Voraussetzung anzunehmen, dass der sechste Teil aller Witwen sich wieder verheiratet und ihre Pensionen damit erlöschen; \textit{oder} er gibt Zahlen an, die erheblich \textit{mehr} Witwen ergeben als die halbe Zahl der bestehenden Ehen; zum Beispiel:\par}

\begin{quote}
in der eben angeführten Schrift S. 27 wird eine Rechnung angestellt, nach der auf 3293 Ehepaare 1831 Witwen kommen,

Prüfung einer kleinen Schrift u. s. w. S. 17, wo auf 100 Ehemänner 60 Witwen gerechnet werden.
\end{quote}
\noindent Eben dieses Verhältnis ist angegeben in
\begin{quote}
Sammlung dreier Aufsätze über die Calenbergische, Preussische und Dänische Witwenversorgungsanstalten S. 30—35.

Sammlung wichtiger Erfahrungen u. s. w. S. 35, wo \textsc{Kritter} auf 450 Ehemänner 280 Witwen rechnet.
\end{quote}
\noindent\hspace*{\parindent}Im Übrigen hat \textsc{Kritter} seine Angaben nirgends auf wirkliche unmittelbare Erfahrungen gestützt (solche Erfahrungen von genügender Art waren damals aus Deutschland auch schwerlich zu beschaffen), sondern auf die Sterblichkeitstafeln und auf gewisse Voraussetzungen über die durchschnittlichen Altersverhältnisse der in die Gesellschaft Eintretenden. Diese Voraussetzungen sind jedenfalls nicht ungünstiger gewählt, als die Verhältnisse bei unserer Professorenwitwenkasse tatsächlich sind.

{\rightskip=0pt plus 1em
Wirkliche Erfahrungen aus einem großen und mit unserer Witwenkasse wohl ungefähr vergleichbaren Kreis finde ich in \textsc{Price}, Observations on reversionary payments, S. 79. 269. 276 (nach der dritten Ausgabe dieses Werks von 1773). Dort wird für die Gesamtheit der Pfarrer und Professoren in Schottland im Durchschnitt von 17 Jahren die Zahl der bestehenden Ehen mit 667 und die der Witwen mit 380 angegeben, also sehr nahe im Verhältnis von 7 zu 4. Zugleich wird bemerkt, dass diese Berufsgruppen dort durchschnittlich im Alter von 27 Jahren in den Genuss des Einkommens aus ihren Stellen kamen, also gewiss früher, als es durchschnittlich bei den Professoren deutscher Universitäten angenommen werden darf. Auch waren bei jenen Witwen die Wiederverheiratungen nicht so selten wie bei den Witwen Göttinger Professoren.\par}


Aus den eigenen unmittelbaren Erfahrungen unserer Gesellschaft lässt sich das für sie gültige Normalverhältnis, auch abgesehen von ihrem geringen Umfang, schon deshalb nicht ableiten, weil die Entwicklung der Anzahl der in ihr bestehenden Ehen eine unbekannte Größe ist. Diese Anzahl ist nämlich aus der gesamten hundertjährigen Dauer der Einrichtung nur für zwei Zeitpunkte bekannt: für den gegenwärtigen Augenblick und nach P.s Angabe (oben S. 132) für den Zeitpunkt der Verhandlungen
\SourceLeaf{151}
von 1794. Hätte man für sämtliche 204 Professoren, die der Witwenkasse bisher beigetreten sind, regelmäßig aufgezeichnet, ob und während welchen Teils ihrer Mitgliedschaft sie verheiratet waren, zugleich mit der genauen Altersangabe für sie selbst und ihre Frauen, so könnte dies mithilfe der Sterblichkeitstafeln zu einer sehr genauen indirekten Bestimmung des fraglichen Verhältnisses genutzt werden. Von all dem ist aber meines Wissens nichts geschehen. Je mehr dies jetzt beklagt werden muss, desto zuversichtlicher darf wohl gehofft werden, dass geeignete Maßnahmen, die wenigstens von jetzt an zu treffen sind, demjenigen, der nach weiteren 100 Jahren ein Gutachten abgeben wird, dieselbe Klage ersparen werden.

